% Section 3.3, from lectures/L03/examples/cpp_interface: the README's prose, then every file of
% the example typeset straight from the repository.

\section{Example: An LED Interface in C++}
\label{c3:sec:example}
This is a demonstration of an interface \code{driver::led::Interface}, which is used to easily
work with different types of LEDs, for example LEDs from different processors, via subclasses. It
is the complete program behind \secref{c3:sec:usinginterfaces}, and it lives in
\file{lectures/L03/examples/cpp_interface}.

\subsection{Files}
\label{c3:sec:examplefiles}
\begin{itemize}
  \item \file{driver/led/interface.hpp} constitutes the interface itself in the form of the base
    class \code{driver::led::Interface}.
  \item \file{driver/led/atmega328p.hpp} constitutes an implementation of LEDs for the processor
    ATmega328P via the subclass \code{driver::led::Atmega328p}. The corresponding source code is
    \file{source/driver/led/atmega328p.cpp}.
  \item \file{driver/led/esp32s3.hpp} constitutes an implementation of LEDs for the processor
    ESP32-S3 via the subclass \code{driver::led::Esp32s3}. The corresponding source code is
    \file{source/driver/led/esp32s3.cpp}.
  \item In \file{main.cpp}, LEDs for each processor are created (via the previously mentioned
    subclasses) and blink at different frequencies:
  \begin{itemize}
    \item This is achieved through a shared blink function that expects a reference to an LED
      interface (\code{driver::led::Interface}).
    \item Since both subclasses inherit \code{driver::led::Interface}, this works perfectly.
  \end{itemize}
\end{itemize}

\subsubsection{The interface}
\cppfile{lectures/L03/examples/cpp_interface/include/driver/led/interface.hpp}

\subsubsection{The ATmega328P implementation}
\cppfile{lectures/L03/examples/cpp_interface/include/driver/led/atmega328p.hpp}
\cppfile{lectures/L03/examples/cpp_interface/source/driver/led/atmega328p.cpp}

\subsubsection{The ESP32-S3 implementation}
\cppfile{lectures/L03/examples/cpp_interface/include/driver/led/esp32s3.hpp}
\cppfile{lectures/L03/examples/cpp_interface/source/driver/led/esp32s3.cpp}

\subsubsection{The program}
\cppfile{lectures/L03/examples/cpp_interface/source/main.cpp}

\subsection{Compiling and running the program}
\label{c3:sec:examplebuild}
The example is built with the Makefile from \secref{c1:sec:tips}, extended with the three source
files and the \file{include} directory:

\makefile{lectures/L03/examples/cpp_interface/makefile}

\noindent Compile and run the program with the following command (in this directory):

\begin{shell}
make
\end{shell}

\noindent You can also build the program without running it afterwards with the following command:

\begin{shell}
make build
\end{shell}

\noindent You can run the program without compiling beforehand with the following command:

\begin{shell}
make run
\end{shell}

\noindent You can also remove compiled files with the following command:

\begin{shell}
make clean
\end{shell}
